\chapter{Mortar Methods and Nonconforming Mesh Coupling}\label{ch:13}
\begin{chaptersummary}
Local high-order spaces allow independent element sizes and degrees, so interface nodes and trace spaces may not match. Mortars organize projection, flux evaluation, and return transfer on a shared interface, ensuring that both sides use the same integral data. Projection matrices, constant preservation, and linear energy conditions provide an interface for hp adaptation and heterogeneous methods.
\end{chaptersummary}
\learninggoals{Construct a common interface partition and mortar space; calculate trace projection and dual return transfer; prove cancellation of integrated fluxes; state the quadrature and trace-reproduction conditions for energy stability.}
\studyroute{Inconsistent curved geometry, nonlinear entropy projections, and specialized reduced-order mortar stabilization are omitted. The linear proof uses a common space reproducing both traces.}{mortar method; nonconforming interface; mass-adjoint projection; hp nonconforming DG}{On one segment, construct mortar mass and projection matrices for a linear left trace and constant right trace, then check flux loads and constant tests explicitly.}

\section{Common Interfaces and Trace Spaces}\index{Mortar methods}
Consider a planar edge $\Gamma$ between two two-dimensional elements, with normal fixed from left to right.
For unequal polynomial degrees, or one coarse edge adjoining several fine edges, first partition the physical interface into common subfaces.
On each subface, choose $V_m=\Poly_{p_m}$ and use identical geometry mappings, normals, and integration measures.
Taking $p_m$ at least as large as both trace degrees reproduces both traces on affine subfaces.
For this conservative coupling approach, see~\cite{kopriva1996}.

\conceptbridge{Why introduce a common interface instead of evaluating fluxes at each side's own nodes?}
{Different nodes lead to different flux samples and integrals. A common interface represents both traces in one space,
evaluates one numerical flux, and returns it with opposite outward normals. Constant preservation ensures agreement of total quantities;
compatible integration and dual transfer also connect the two energy pairings.}
{Nonmatching traces $\longrightarrow$ common interface space $\longrightarrow$ shared flux and compatible return transfer.}

\section{Projection, Fluxes, and Dual Return Transfer}
Let $\mu_\alpha$ be a mortar basis and $\phi_i^s$ a trace basis on side $s=L,R$. Define
\[
(M_m)_{\alpha\beta}=\int_\Gamma\mu_\alpha\mu_\beta\dd S,\qquad
(B_{ms})_{\alpha i}=\int_\Gamma\mu_\alpha\phi_i^s\dd S.
\]
The $L^2$ trace projection satisfies
\begin{equation}\label{eq:mortar-proj}
P_s=M_m^{-1}B_{ms},\qquad u_{s,m}=P_su_s.
\end{equation}
Evaluate a single-valued normal flux from $u_{L,m},u_{R,m}$ at common quadrature nodes.
Represent its load with the same rule as $M_mf_m^*$. The outward flux loads are
\begin{equation}\label{eq:mortar-return}
r_L=P_L^{\mathsf T}M_mf_m^*,\qquad
r_R=-P_R^{\mathsf T}M_mf_m^*.
\end{equation}
These are dual loads in weak residuals. Flux coefficients on one side require an additional inverse trace mass matrix;
this also gives the mass-adjoint relation $M_sP_{m\to s}=P_s^{\mathsf T}M_m$.

\section{Conservation and Linear Energy Conditions}
Let $\bm1_s$ and $\bm1_m$ represent constants. Constant reproduction means
$P_s\bm1_s=\bm1_m$, hence
\[
\bm1_L^{\mathsf T}r_L+\bm1_R^{\mathsf T}r_R=0.
\]
If one coarse edge meets several mortar subfaces, sum their loads;
the coarse and fine sides still use the same integrals.

For a linear energy estimate, further assume that the common space reproduces both traces exactly
and that quadrature exactly integrates the relevant products. Then
$P_s^{\mathsf T}M_mP_s=M_s$,
For multiple subfaces, sum the matrices on the left.
For constant normal speed $a_n$, use
\[
f_m^*=a_n\{u_m\}-\frac{|a_n|}{2}[u_m],
\]
The combined DG interface energy contribution is
\[
-\frac{|a_n|}{2}[u_m]^{\mathsf T}M_m[u_m]\le0.
\]
If the common space has lower degree or the geometry and quadrature differ, the trace-norm identity requires a corresponding replacement construction.
Constant preservation supplies total conservation; trace reproduction, quadrature, and mass-adjoint transfer additionally support the energy estimate.

\section{A Hand-Computable Nonconforming Interface}
Let $\Gamma=[-1,1]$, with left trace basis $(1,s)$ and a constant-only right trace.
Use mortar basis $(1,s)$. Exact integration gives
\[
M_m=\begin{pmatrix}2&0\\0&2/3\end{pmatrix},\quad
P_L=I_2,\quad P_R=\begin{pmatrix}1\\0\end{pmatrix}.
\]
For shared flux coefficients $(F_0,F_1)^{\mathsf T}$,
\[
r_L=\begin{pmatrix}2F_0\\2F_1/3\end{pmatrix},\qquad r_R=-2F_0.
\]
a constant test reads only the first left component, so the total flux cancels exactly.
The linear component contributes to the left trace residual. Changing interface spaces during hp adaptation requires updating projections and physical measures together.

\aiexercise{Use a quadratic left trace and linear right trace, with three Gauss points for exact integration. Construct $M_m,B_{mL},B_{mR}$ and verify constant reproduction, mass adjointness, and upwind interface energy.}
{Split the right edge into two subfaces with linear traces. Write the subface-map Jacobians and prove that the coarse-side total load is the negative sum of the two fine-side loads.}
